% Generated draft: requires mathematical and source review.
\documentclass[11pt]{article}
\usepackage[margin=25mm]{geometry}
\usepackage{fontspec}
\setmainfont{Latin Modern Roman}
\usepackage{amsmath,amssymb,mathtools,unicode-math}
\setmathfont{Latin Modern Math}
\usepackage{enumitem,needspace,xurl,hyperref}
\setlength{\parindent}{0pt}
\setlength{\parskip}{5pt}
\setlength{\emergencystretch}{3em}
\newcommand{\N}{\mathbb{N}}
\newcommand{\Z}{\mathbb{Z}}
\newcommand{\Q}{\mathbb{Q}}
\newcommand{\R}{\mathbb{R}}
\newcommand{\C}{\mathbb{C}}
\newcommand{\F}{\mathbb{F}}
\newcommand{\A}{\mathbb{A}}
\newcommand{\PP}{\mathbb{P}}
\newcommand{\E}{\mathbb{E}}
\newcommand{\eps}{\varepsilon}
\newcommand{\dd}{\,\mathrm{d}}
\begin{document}
\section*{UnsolvedMath ID: 30002109}
\textbf{Source problem code:} catalogue entry
\subsection{Definitions and mathematical statement}
Let \(a\) denote the separation parameter of the two domains in the supplied plasmonic boundary-value problem. The singular coalescence limit is \(a\to0\), taken through admissible parameter values. For notation only, write \(D_1(a)\) and \(D_2(a)\) for these domains and \(\mathcal P_a\) for the associated boundary-value problem. These symbols do not specify an ambient space, geometry, differential operator, material parameters, boundary or transmission conditions, or solution space; none of these is given in the supplied record.

\textbf{Questions in the singular limit.} For the coalescing configuration described by the source, determine:
\begin{enumerate}
\item whether solutions exist at \(a=0\);
\item what modification of the equations is required to formulate the problem at \(a=0\);
\item whether qualitatively different solution branches occur, including branches described in the source as localized or hybridized.
\end{enumerate}
The limiting problem \(\mathcal P_0\) is itself part of what is to be determined, rather than a prescribed problem obtained by formally substituting \(a=0\). The source does not define localization, hybridization, or a solution branch. It also does not specify whether existence at \(a=0\) means independent solvability of a limiting problem, convergence of solutions for nonzero \(a\), or both. Consequently, these questions faithfully retain the supplied research target, but a fully quantified boundary-value formulation cannot be recovered from this record alone.
\subsection{Short English statement}
As the two domains in a plasmonic boundary-value problem coalesce with \(a\to0\), determine whether solutions exist at zero separation, which limiting equations apply, and whether localized or hybridized solution branches occur. The supplied record omits the governing equations, geometry, solution class, and definitions of these branches, so a precise analytic formulation requires further source information.
\subsection{Sources}
\begin{itemize}
\item Local UnsolvedMath registry: supplied complete record with permanent ID \(30002109\), source problem code \texttt{catalogue entry}, in the Oberwolfach Reports Open Problems collection. No individual proposer is identified; the supplied proposed year is 2012.
\item Supplied report citation: \emph{Geometric Aspects of Spectral Theory (2013)}, Oberwolfach Reports, pp.~2013--2076, DOI \texttt{10.4171/owr/2012/33}. \url{https://doi.org/10.4171/owr/2012/33}. The record attributes the existence, limiting-formulation, and branch questions to this report. Its \texttt{corrected\_verified} label and statement-verification assertion are supplied registry claims, not independent verification here.
\item Supplied literature source: Helsing and Perfekt, \emph{the plasmonic eigenvalue problem}. \url{https://arxiv.org/abs/1208.3120}.
\item Supplied literature source: \emph{geometric quantization of localized surface plasmons}. \url{https://academic.oup.com/imamat/article-abstract/84/4/813/5536913}.
\item Supplied literature source: original OWR report. \url{https://doi.org/10.4171/owr/2012/33}.
\item Literature assessment dated exactly \texttt{2026-08-21}: the supplied curation labels the problem partially solved and states that plasmonic eigenvalue asymptotics and blow-up methods resolve important smooth and nearly singular geometries, while the general singular limiting equations and all qualitatively different branches requested in the record remain unsettled. This is supplied dated curation derived from the OWR material, not an independently verified literature assessment; its status and cited claims require independent checking.
\end{itemize}
\end{document}
